\section{优势与代价：为什么 Transformer 成功，也为什么仍难}

Transformer 的成功通常归因于长距离依赖和并行训练，但这两个优势都带条件。Full attention 提供短信息路径，却创建 $n^2$ score；矩阵计算适合 GPU/TPU，却要求批量、内存和数据规模；较弱的局部归纳偏置提升通用性，也可能增加样本需求。

\subsection{复杂度与系统边界}

单层 full self-attention 的 score 计算约为 $O(n^2d)$，score memory 约为 $O(n^2)$；RNN 单层通常约 $O(nd^2)$ 并需要 $O(n)$ sequential steps。哪一种更快取决于 $n$、$d$、batch、kernel、memory bandwidth 和硬件。复杂度表不能直接等于 wall-clock latency。

\lecturefigure{slide-15-good-bad.jpg}{Transformer 的长程建模与并行性伴随二次 attention、数据需求和较弱局部偏置。}{官方视频 16:10。}

\begin{knowledgebox}{读图：优点与缺点来自同一设计}
任意 token 直接交互带来长程依赖，也产生 dense pairwise matrix；不依赖 recurrence 带来训练并行，也减少内建顺序偏置；统一 block 方便扩展，也可能需要更多数据学习领域结构。读图时不要把左右两列当独立清单，而要寻找同一设计选择的两面。
\end{knowledgebox}

\begin{warningbox}{FlashAttention 不会把数学复杂度变成线性}
Fused/tiled attention 可以减少 \term{HBM}（High Bandwidth Memory，GPU 高带宽显存）读写和中间矩阵物化，显著改善实际速度与内存，但标准 full attention 的 pairwise 计算仍是二次。Algorithmic sparsity、low-rank、state-space 或 recurrence 是不同路线；系统优化与算法复杂度必须分开描述。
\end{warningbox}

\begin{knowledgebox}{术语消化：inductive bias 与 sample efficiency}
\term{inductive bias}（归纳偏置）是模型在看数据前就偏好的结构，例如 convolution 的局部平移共享、RNN 的时间递归、causal mask 的方向性。较少偏置提供更大通用性，但模型可能需要更多数据才能学会本可手工编码的规律。\term{sample efficiency} 描述达到目标性能所需数据，而不是每秒处理样本数。
\end{knowledgebox}

\subsection{本章小结}

Transformer 的优势与代价同源：动态全局连接、矩阵并行和弱领域假设扩大了可扩展性，也带来二次成本、数据需求与系统复杂度。工程判断必须同时看算法和硬件。

